\begin{block}{Preliminary results I}
  Basic array operations (within a chunk) perform better in the specialized runtime.

  \begin{figure}
    \includegraphics[height=20cm]{place_holder.png}

    Box-and-whisker graph showing the performance (measured in cycles) and uncertainty
    \caption{Nautilus can run array operations faster than Linux.}
  \end{figure}

  % \begin{table}
  %   \begin{tabular}{l || c c | c c}
  %     Operation & Linux & & Nautilus & \\
  %     \hline\hline
  %     Allocate & 1.0e9 & 1.0x & 1.0e9 & 1.0x \\\hline
  %     Free & 1.0e9 & 1.0x & 1.0e9 & 1.0x \\\hline
  %     Copy (\texttt{memcpy}) & 1.0e9 & 1.0x & 1.0e9 & 1.0x \\\hline
  %     Copy (individually) & 1.0e9 & 1.0x & 1.0e9 & 1.0x \\\hline
  %     Reverse & 1.0e9 & 1.0x & 1.0e9 & 1.0x \\
  %   \end{tabular}

  %   Runtime measured in cycles according to \texttt{rdtsc}.

  %   \caption{Performance of Linux vs Nautilus on basic array operations averaged over 10 runs.}
  % \end{table}
  
\end{block}

\begin{block}{Preliminary results I (determinism)}

  \begin{figure}
    \includegraphics[height=20cm]{place_holder.png}

    \todo{Line plot with $\pm \sigma$ region shaded.}
  \end{figure}

  \begin{table}
    \begin{tabular}{l || c | c }
      Operation & Linux & Nautilus \\
      \hline\hline
      TLB miss & & 1 \\
      Page fault & & 0 \\
      Context switches & & 0 \\
      Interrupts & & 0 \\
    \end{tabular}

    \todo{measure linux with perf.}
    \caption{Nautilus can run array operations faster than Linux.}
  \end{table}

\end{block}
%%% Local Variables:
%%% mode: latex
%%% TeX-master: "main"
%%% End:
